\chapter{Discusión}

La discusión no repite el inventario del Capítulo~\ref{sec:resultados-condiciones}; relaciona cada bloque de evidencia con la afirmación científica que puede sostener. Mientras la campaña final permanezca pendiente, los razonamientos que dependen de ella se expresan como criterios de decisión y no como resultados anticipados.

\section{Rendimiento intradistribución frente a generalización}

Una partición aleatoria estratificada responde a una pregunta necesaria pero limitada: si entrenamiento y prueba se extraen de la misma mezcla de CICIDS2017, ¿puede el pipeline aprender la regla binaria definida? La MAIN histórica responde afirmativamente para una ejecución anterior. Sin embargo, esta pregunta no equivale a saber si el modelo reconocerá ataques de otro día, captura o distribución. Los patrones de tráfico, herramientas, prevalencias y familias de ataque varían entre los CSV del benchmark, y una mezcla aleatoria diluye esa separación.

La campaña final trata la partición aleatoria como control intradistribución y reserva la afirmación de generalización al contraste con el experimento completo por días y los escenarios retenidos. Si MAIN mantiene métricas altas pero estas caen al separar días o CSV, la interpretación correcta será que el modelo explota bien la distribución mezclada, no que ofrece una defensa robusta. Si el comportamiento persiste, la afirmación podrá reforzarse únicamente para esos desplazamientos observados; seguirá sin equivaler a producción independiente.

\section{Duplicados y continuidad entre particiones}

El análisis histórico encontró una proporción relevante de coincidencias exactas entre entrenamiento y prueba. Esto no prueba una fuga de etiqueta dentro del código ni invalida automáticamente toda evaluación aleatoria: demuestra que el conjunto de prueba contiene menos novedad de la que su tamaño sugiere. En consecuencia, los intervalos estrechos obtenidos por \emph{bootstrap} pueden coexistir con una validez externa débil. El primero mide precisión sobre un test fijo; la segunda depende de la separación que ese test proporciona.

La regeneración final del análisis de duplicados debe interpretarse junto con el día completo. Si la tasa de coincidencias vuelve a ser elevada y el rendimiento temporal disminuye, ambas señales ofrecerán una explicación coherente del optimismo aleatorio. Si el rendimiento temporal se conserva, los duplicados seguirán siendo una amenaza del benchmark, pero no bastarán para atribuirles el resultado del modelo.

\section{Lectura sensible al coste}

La recompensa distingue cuatro consecuencias: permitir tráfico benigno, bloquear un ataque, bloquear tráfico benigno y permitir un ataque. La penalización más severa del falso negativo expresa una preferencia defensiva, pero no elimina el coste del falso positivo. Por eso la exactitud agregada no basta: deben leerse conjuntamente \emph{recall} de ataque, tasa de falsos negativos, precisión de ataque, tasa de falsos positivos y matriz de confusión.

Esta lectura es especialmente importante fuera de la partición aleatoria. Un modelo puede conservar una precisión alta de ataque porque bloquea muy pocos flujos y, a la vez, omitir la mayoría de las intrusiones. También puede mejorar \emph{recall} a costa de interrumpir tráfico legítimo. La campaña no fijará un umbral operativo ni ejecutará bloqueos reales; comparará el compromiso producido por las políticas entrenadas bajo una función de recompensa explícita.

\section{QRDQN frente a Random Forest}

Random Forest no se incluye como adversario débil, sino como prueba de necesidad. Con $\gamma=0$ y filas independientes, QRDQN resuelve una tarea próxima a clasificación tabular sensible al coste. Un clasificador supervisado sólido bajo el mismo esquema, partición y métrica permite preguntar qué aporta realmente la formulación distribucional, en lugar de atribuir a RL cualquier rendimiento que provenga de las características o del benchmark.

La evidencia histórica muestra que RF puede dominar la partición aleatoria y fallar ante una separación temporal concreta. Ese contraste es indicativo, pero está desbalanceado porque el QRDQN temporal histórico era un proxy. La campaña final corrige esta asimetría con pares equivalentes en MAIN, un millón de filas, día completo y cuatro escenarios retenidos. Si RF y QRDQN se degradan de forma parecida, la limitación principal apuntará a datos y dominio. Si difieren de forma consistente, podrá discutirse una sensibilidad específica de la familia de modelo. Un resultado favorable aislado no justificará superioridad general.

\section{Conjunto de datos como entorno y $\gamma=0$}

El uso de un conjunto tabular como entorno aporta una interfaz explícita de acciones y recompensas, compatibilidad con instrumentación de RL y una base para estudiar políticas sensibles al coste. A la vez, con $\gamma=0$ la devolución no incorpora consecuencias futuras: cada decisión depende de la observación actual y de su recompensa inmediata. La distribución cuantílica modelada por QRDQN describe la devolución de esa decisión, no una secuencia causal de estados de red.

Esta formulación enseña una lección metodológica: llamar \enquote{entorno} a un iterador sobre filas no convierte automáticamente la detección en un problema secuencial. La comparación supervisada es, por tanto, parte constitutiva de la evaluación. Una extensión genuinamente secuencial exigiría definir estado temporal, consecuencias de las acciones, episodios con dinámica de red y riesgos de intervención; queda fuera de este trabajo.

\section{Eficiencia de datos y sensibilidad a la semilla}

La escalera de tamaños distinguirá dos explicaciones que una única MAIN no puede separar: que el comportamiento dependa de casi todo CICIDS2017 o que se alcance con menos datos. Los presupuestos de pasos crecen proporcionalmente hasta el perfil completo, por lo que la curva debe leerse como respuesta conjunta a volumen y presupuesto de optimización. No medirá una ley universal de escalado.

El estudio de semillas responde a otra incertidumbre. Mantiene $1\,000\,000$ de filas y la semilla de partición~42, y cambia solo la semilla del modelo entre 42 y 46. Así estima sensibilidad de inicialización y aprendizaje sobre una partición fija; no incluye variación del conjunto de prueba. Media, dispersión, rango y casos atípicos deberán reportarse juntos. Una diferencia pequeña entre cinco ejecuciones reforzará la estabilidad local de ese punto, pero no sustituirá más particiones ni el MAIN completo repetido.

\section{Sensibilidad por escenario}

Los cuatro \emph{holdouts} ---Web Attacks, Infiltration, PortScan y DDoS--- fueron elegidos como escenarios dirigidos, no como barrido exhaustivo de los ocho CSV. La comparación debe mostrar el soporte y prevalencia de cada test junto con las métricas: un F1 no es comparable si la composición o la clase positiva cambian sustancialmente. El patrón entre escenarios será más informativo que una media única, porque permitirá identificar qué distribución retenida provoca cada error.

Esta selección limita la inferencia. Incluso si los cuatro escenarios fueran favorables, no quedaría demostrada invariancia frente a cualquier CSV, ataque o día. Si alguno falla, tampoco se deduce que la arquitectura sea inútil; sí se identifica una frontera concreta que requiere análisis de características, cobertura y posible cambio de dominio.

\section{Fase 2 y cambio de dominio}

La Fase~2 comprueba que el modelo, el escalador, los percentiles y el contrato canónico pueden aplicarse a flujos externos al cargador de CICIDS2017. Sus diagnósticos ---presencia de características, valores no finitos, recorte y puntuaciones estandarizadas--- ayudan a distinguir un error del pipeline de un desplazamiento de entrada. Cuando existan etiquetas, la matriz de confusión añade una señal de comportamiento.

No obstante, la captura pertenece a un único laboratorio cerrado, fue generada por el operador y no representa adversarios ni operación independiente. Un resultado alto demostraría compatibilidad y comportamiento sobre esa captura; uno bajo señalaría una limitación de transferencia o del mapeo. Ninguno de los dos, por sí solo, probaría preparación para despliegue.

\section{Amenazas a la validez y evidencia más fuerte}

La validez interna depende del contrato de partición, escalado solo con entrenamiento, semillas explícitas, control de etiquetas barajadas, validación directa y artefactos sellados. La validez de constructo está limitada por el mapeo binario y por tratar decisiones independientes como bandido contextual. La validez externa está limitada por CICIDS2017, cuatro escenarios seleccionados y una captura de laboratorio. La validez de conclusión estará condicionada por cinco semillas en un único tamaño y por la ausencia de repetición completa del MAIN.

Evidencia más fuerte requeriría capturas etiquetadas e independientes de varios entornos, repetición de los principales contrastes, comparación con más familias tabulares bajo el mismo contrato y evaluación previa a cualquier acción real. Esta tesis no pretende cubrir ese último nivel. Su objetivo es dejar una cadena reproducible y honesta que permita saber qué afirmación respalda cada experimento y cuál sigue fuera de alcance.
